\documentclass[11pt,leqno]{amsart}
\usepackage[letterpaper,margin=1in]{geometry}
\usepackage[T1]{fontenc}
\usepackage{lmodern}
\usepackage{amsmath,amssymb,amsthm,mathtools}
\usepackage{microtype}
\usepackage{needspace}
\usepackage{xurl}
\usepackage[colorlinks=true,linkcolor=blue,citecolor=blue,urlcolor=blue]{hyperref}

\title[The minimum squared Lagrange interpolation integral]{An exact algebraic construction for\newline Erd\H{o}s's interpolation minimum}
\author[W. Howe]{Will Howe}
\author[G. Huffman]{Grant Huffman}
\author[J. Kohlmeyer]{Jensen Kohlmeyer}
\author[L. Kruer]{Liam Kruer}
\thanks{\raggedright\textit{E-mail addresses.}
\mbox{\href{mailto:howew@purdue.edu}{\texttt{howew@purdue.edu}} (W.~Howe)},
\mbox{\href{mailto:huffma97@purdue.edu}{\texttt{huffma97@purdue.edu}} (G.~Huffman)},
\mbox{\href{mailto:jenwin@purdue.edu}{\texttt{jenwin@purdue.edu}} (J.~Kohlmeyer)},
\mbox{\href{mailto:kruerl@purdue.edu}{\texttt{kruerl@purdue.edu}} (L.~Kruer)}.}
\date{November 11, 2026}
\subjclass[2020]{Primary 41A05; Secondary 13P10, 14P10, 15A63}
\keywords{Lagrange interpolation, extremal nodes, exact algebraic values,
Hermite matrices, trace forms, real root counting}

\newtheorem{theorem}{Theorem}[section]
\newtheorem{proposition}[theorem]{Proposition}
\newtheorem{lemma}[theorem]{Lemma}
\newtheorem{corollary}[theorem]{Corollary}
\theoremstyle{definition}
\newtheorem{definition}[theorem]{Definition}
\theoremstyle{remark}
\newtheorem{remark}[theorem]{Remark}
\numberwithin{equation}{section}

\newcommand{\R}{\mathbb R}
\newcommand{\Q}{\mathbb Q}
\newcommand{\C}{\mathbb C}
\newcommand{\N}{\mathbb N}
\newcommand{\B}{\mathcal B}
\newcommand{\Lz}{\mathcal L_z}
\newcommand{\NF}{\operatorname{NF}}
\newcommand{\sig}{\operatorname{sig}}
\newcommand{\tr}{\operatorname{tr}}
\newcommand{\diag}{\operatorname{diag}}
\newcommand{\abs}[1]{\lvert #1\rvert}
\newcommand{\norm}[1]{\lVert #1\rVert}

\hypersetup{
 pdftitle={An exact algebraic construction for Erdos's interpolation minimum},
 pdfauthor={Will Howe, Grant Huffman, Jensen Kohlmeyer, Liam Kruer},
 pdfsubject={An exact algebraic value rule for the minimum squared Lagrange interpolation integral},
 pdfkeywords={Lagrange interpolation, Erdos problem 1131, Hermite matrix, trace form}
}

\begin{document}
\begin{abstract}
For distinct nodes in $[-1,1]$, consider the integral of the sum of the
squares of their Lagrange cardinal polynomials. We give a finite algebraic
construction of its minimum for every positive number of nodes. For
$n\geq2$, put $k=n-1$ and $d=2n(n-1)$. A rational symmetric polynomial
matrix of order $(d+1)^k$ determines a nonzero rational polynomial of degree
at most $kd(d+1)^k$. The minimum is a specified real root of this polynomial;
its index is obtained from signatures at rational sample points. The proof
establishes attainment and endpoint inclusion, expresses the objective in
squared gap coordinates, and converts its threshold into a real critical
point count. A fixed monomial basis and the Hermite trace form give the
required count at every parameter, including parameters with multiple
critical points. All steps of the value rule use rational polynomial
arithmetic, univariate real root isolation, and rational symmetric elimination.
\end{abstract}
\maketitle

\section{Introduction and main result}

Let $X=(x_1,\ldots,x_n)$ consist of distinct real nodes in $[-1,1]$.
Its Lagrange cardinal polynomials are
\begin{equation}\label{eq:cardinals}
 \ell_i^X(x)=\prod_{j\ne i}\frac{x-x_j}{x_i-x_j},
 \qquad 1\leq i\leq n.
\end{equation}
We use the empty product convention when $n=1$. Define
\begin{equation}\label{eq:objective}
 I(X)=\int_{-1}^{1}\sum_{i=1}^n\ell_i^X(x)^2\,dx,
 \qquad
 m_n=\inf_{\substack{x_i\in[-1,1]\\x_i\ne x_j\ (i\ne j)}}I(X).
\end{equation}
The integral is taken with respect to ordinary Lebesgue measure.
The value of $I(X)$ is invariant under a permutation of the nodes, so
ordered nodes may be used throughout.

This minimization is among the classical interpolation problems of
Erd\H{o}s; see Erd\H{o}s, Szabados, Varma, and V\'ertesi~\cite{ESVV}.
Brutman and Toledano~\cite{BT} studied the finite minima numerically and
their asymptotic behavior. The problem is also recorded as Erd\H{o}s
Problem~1131~\cite{Bloom}. Our result concerns the exact finite-$n$ minimum.

Algebraicity and an effective value procedure follow from general real
algebraic methods~\cite{BPR}. Indeed, after clearing the positive squared
Vandermonde denominator on the ordered node region, the assertion that
$t$ is an attained minimum is a first-order formula over $\Q$.
Quantifier elimination then describes the singleton $\{m_n\}$ by
univariate polynomial sign conditions. The construction below gives
explicit polynomial and matrix data for this particular objective.
It uses the classical Hermite trace principle and its parametric form;
see~\cite{BPR,LeSafey}. We include the trace argument needed here, including
the treatment of multiplicities.

For a real symmetric matrix $H$, write
\[
 \sig H=n_+(H)-n_-(H),
\]
where $n_+(H)$ and $n_-(H)$ count positive and negative eigenvalues with
multiplicity. Zero eigenvalues contribute zero.

\begin{theorem}\label{thm:main}
For every positive integer $n$, the infimum in~\eqref{eq:objective} is
attained. For $n=1$ it equals $2$. For $n\geq2$, every minimizing node set
contains both endpoints, and
\begin{equation}\label{eq:range}
 \frac{2}{n}\leq m_n<2.
\end{equation}
Fix $n\geq2$ and put
\[
 k=n-1,\qquad d=2n(n-1),\qquad M=(d+1)^k.
\]
There are an explicitly defined symmetric polynomial matrix and a
nonzero polynomial satisfying
\[
 H_n(t)\in\operatorname{Sym}_{M}(\Q[t]),\qquad
 0\ne R_n(t)\in\Q[t],\qquad \deg R_n\leq kdM,
\]
constructed in Sections~\ref{sec:quotient} and~\ref{sec:selection}, with
the following property. Let $r_1<\cdots<r_s$ be the distinct real roots
of $R_n$ in $(0,2)$, put $r_{s+1}=2$, and choose rational numbers
$\rho_j\in(r_j,r_{j+1})$ for $1\leq j\leq s$. Then $s\geq1$ and
\begin{equation}\label{eq:main-rule}
 \boxed{\displaystyle
 m_n=r_{1+c_n},\qquad
 c_n=\#\{j:1\leq j\leq s,\ \sig H_n(\rho_j)=1\}.}
\end{equation}
The construction and the rule depend only on $n$ and use finitely many
exact rational operations and univariate real root isolations.
\end{theorem}

The proof proceeds by an endpoint reduction, a homogeneous polynomial
encoding, and a coercive perturbation whose real critical points detect
the minimum. A fixed quotient basis makes the trace matrix valid at
every real parameter. The last step identifies the correct root by
counting signatures on the intervals between consecutive roots.

\section{Attainment and endpoint reduction}\label{sec:endpoints}

Write $S_X(x)=\sum_i\ell_i^X(x)^2$.

\begin{lemma}\label{lem:attainment}
The minimum $m_n$ is attained, and $m_n\geq2/n$.
\end{lemma}

\begin{proof}
Choose a basis $p_0,\ldots,p_{n-1}$ of the polynomials of degree at most
$n-1$ that is orthonormal for the inner product
$\langle f,g\rangle=\int_{-1}^1fg$. Let
$V_{ij}=p_{j-1}(x_i)$. Distinctness of the nodes makes $V$ invertible.
The coefficient vector of $\ell_i^X$ in this basis is $V^{-1}e_i$.
Consequently,
\begin{equation}\label{eq:gram}
 I(X)=\norm{V^{-1}}_F^2
      =\tr\bigl((V^{\mathsf T}V)^{-1}\bigr).
\end{equation}
Here $\norm{\cdot}_F$ is the Frobenius norm.

Suppose a sequence of ordered configurations tends to a collision in
the closed ordered region
$-1\leq x_1\leq\cdots\leq x_n\leq1$. The limiting evaluation matrix has
two equal rows, so its smallest singular value tends to zero. Since
\[
 I(X)\geq\sigma_{\min}(V)^{-2},
\]
the objective tends to infinity along such a sequence.
A minimizing sequence has bounded objective by comparison with any
fixed distinct-node configuration. Compactness of the closed ordered
region therefore gives a convergent subsequence with a distinct-node
limit. Continuity of~\eqref{eq:gram} at that limit proves attainment.
For $n=1$ the same conclusion follows directly from $\ell_1^X=1$.

The interpolation identity $\sum_i\ell_i^X(x)=1$ and the Cauchy--Schwarz
inequality give $S_X(x)\geq1/n$. Integration proves the lower bound.
\end{proof}

\Needspace{9\baselineskip}
\begin{lemma}\label{lem:strict-upper}
For every $n\geq2$, there is a node configuration with $I(X)<2$.
\end{lemma}

\begin{proof}
Maximize the Vandermonde product
\[
 \Delta(X)=\prod_{a<b}(x_b-x_a)
\]
on the closed ordered node region. The maximum is positive, hence the
maximizing nodes are distinct. Moving the leftmost node to $-1$ strictly
increases the product unless it is already there. The corresponding
argument at the right gives $x_n=1$.

For every interior maximizing node, logarithmic differentiation yields
\begin{equation}\label{eq:fekete-derivative}
 0=\sum_{j\ne i}\frac{1}{x_i-x_j}=(\ell_i^X)'(x_i).
\end{equation}
The polynomial $S_X$ equals $1$ at each node. At an interior node,
\[
 S_X'(x_i)=2\sum_j\ell_j^X(x_i)(\ell_j^X)'(x_i)
         =2(\ell_i^X)'(x_i)=0.
\]
Moreover, $S_X$ has degree exactly $2n-2$, since its leading coefficient
is a sum of positive squares. Counting the endpoint zeros and the double
interior zeros of $1-S_X$ now gives
\begin{equation}\label{eq:fekete-factor}
 1-S_X(x)=c(1-x^2)\prod_{i=2}^{n-1}(x-x_i)^2,\qquad c>0.
\end{equation}
The sign of $c$ follows from the leading coefficient. The right-hand
side is nonnegative on $[-1,1]$ and is positive except at finitely many
points. Thus $I(X)<2$.
\end{proof}

\begin{proposition}\label{prop:endpoints}
Every minimizing node set with $n\geq2$ contains $-1$ and $1$.
\end{proposition}

\begin{proof}
Let $X$ be a minimizing configuration with hull $[a,b]$, and put
$L=b-a>0$. For $x<a$, every factor of
\[
 \ell_1^X(x)=\prod_{j>1}\frac{x-x_j}{a-x_j}
\]
is greater than one. Thus $S_X(x)>1$ on $[-1,a)$. The same argument with
$\ell_n^X$ gives $S_X(x)>1$ on $(b,1]$.

Suppose $L<2$, and write
\[
 A=\int_a^b S_X(x)\,dx,\qquad
 B=\int_{[-1,1]\setminus[a,b]}S_X(x)\,dx.
\]
Then $B>2-L$. Lemma~\ref{lem:strict-upper} gives $A+B=m_n<2$, so $A<L$.
The affine map from $[a,b]$ onto $[-1,1]$ carries $X$ to an admissible
configuration $Y$. Its cardinal polynomials transform by composition,
and change of variables gives $I(Y)=2A/L$. Consequently,
\[
 I(X)-I(Y)=B-\frac{2-L}{L}A
          >(2-L)\left(1-\frac{A}{L}\right)>0,
\]
contradicting minimality. Hence $L=2$, which forces $a=-1$ and $b=1$.
\end{proof}

The affine hull argument also occurs in the exact four-node analysis
of White~\cite{White}. Proposition~\ref{prop:endpoints} gives the
reduction for every $n\geq2$.

\section{A rational homogeneous encoding}\label{sec:encoding}

Fix $n\geq2$, and retain $k=n-1$ and $d=2n(n-1)$. For
$u=(u_1,\ldots,u_k)$ define
\begin{equation}\label{eq:gaps}
 g_a=u_a^2,\qquad T=\sum_{a=1}^k g_a,\qquad
 Y_i=-T+2\sum_{a<i}g_a\quad(1\leq i\leq n).
\end{equation}
In particular, $Y_1=-T$, $Y_n=T$, and $Y_{i+1}-Y_i=2u_i^2$.
For an $n$-tuple $y$, let $y^{(i,w)}$ denote the tuple obtained by
replacing its $i$th entry by $w$.

Define the $\Q[u]$-linear functional $\Lz$ on $\Q[u,z]$ by
\begin{equation}\label{eq:moment-functional}
 \Lz(z^j)=
 \begin{cases}
  2/(j+1),&j\text{ even},\\
  0,&j\text{ odd}.
 \end{cases}
\end{equation}
Thus $\Lz$ is integration in $z$ over $[-1,1]$. Put
\begin{equation}\label{eq:ND}
 D(u)=\Delta(Y)^2,\qquad
 N(u)=\sum_{i=1}^n\Lz\!\left(\Delta\bigl(Y^{(i,Tz)}\bigr)^2\right).
\end{equation}
These are rational homogeneous forms of degree $d$. Indeed, every
entry of $Y$ and the replacement $Tz$ has $u$-degree two, and the
Vandermonde has degree $n(n-1)/2$ in its entries.

\begin{lemma}\label{lem:encoding}
For real $u$ with every coordinate nonzero,
\begin{equation}\label{eq:ratio}
 \frac{N(u)}{D(u)}=I(Y_1/T,\ldots,Y_n/T).
\end{equation}
Every ordered configuration with endpoints $-1$ and $1$ arises in
this way. For every real $u$, $D(u)\geq0$ and $N(u)\geq0$, and
$D(u)=0$ precisely when some $u_a=0$.
\end{lemma}

\begin{proof}
If all coordinates are nonzero, then $T>0$ and~\eqref{eq:gaps} gives
strictly increasing normalized nodes with the required endpoints.
For any distinct-node tuple $x$, the Vandermonde replacement identity is
\[
 \ell_i^x(z)=\frac{\Delta(x^{(i,z)})}{\Delta(x)}.
\]
Scaling all entries by $T$ shows that this ratio for $x=Y/T$ equals
$\Delta(Y^{(i,Tz)})/\Delta(Y)$. Squaring, summing, and integrating gives
\eqref{eq:ratio}.

Conversely, for ordered nodes with these endpoints choose
\[
 u_a=\sqrt{\frac{x_{a+1}-x_a}{2}}.
\]
Then $T=1$ and $Y_i=x_i$.
Finally, $D$ is a square, and its zero set is determined by the
successive differences $2u_a^2$. Each summand defining $N$ is the
integral of a real square, which proves its nonnegativity even when
several gaps vanish.
\end{proof}

For a real parameter $t$, define
\begin{equation}\label{eq:P}
 P_t(u)=N(u)-tD(u).
\end{equation}

\begin{proposition}\label{prop:threshold}
For every real $t$,
\begin{equation}\label{eq:nonnegative-threshold}
 P_t(u)\geq0\text{ for every }u\in\R^k
 \quad\Longleftrightarrow\quad t\leq m_n.
\end{equation}
\end{proposition}

\begin{proof}
If $t\leq m_n$ and $D(u)>0$, Lemma~\ref{lem:encoding} gives
$P_t(u)=D(u)(I(Y/T)-t)\geq0$. At every remaining real point,
$D(u)=0$ and $P_t(u)=N(u)\geq0$.

Conversely, Lemma~\ref{lem:attainment} and
Proposition~\ref{prop:endpoints} provide a minimizing configuration
with both endpoints. Apply the converse parametrization in
Lemma~\ref{lem:encoding}. Its denominator is positive, so
$P_t(u)\geq0$ implies $m_n\geq t$.
\end{proof}

\section{Detecting the threshold by critical points}\label{sec:critical}

Set
\begin{equation}\label{eq:Q}
 Q_t(u)=P_t(u)+\sum_{a=1}^k u_a^{d+2}.
\end{equation}
Since $d$ is even, the added form is strictly positive away from the origin.

\begin{proposition}\label{prop:critical}
The polynomial $Q_t$ is coercive for every real $t$, and
\begin{equation}\label{eq:critical-threshold}
 \{u\in\R^k:\nabla Q_t(u)=0\}=\{0\}
 \quad\Longleftrightarrow\quad t\leq m_n.
\end{equation}
\end{proposition}

\begin{proof}
On the unit sphere, $\sum_a u_a^{d+2}$ has a positive minimum, while
$\abs{P_t(u)}$ has a finite maximum. Homogeneity therefore gives constants
$c>0$ and $C_t\geq0$ such that
\[
 Q_t(u)\geq c\norm{u}^{d+2}-C_t\norm{u}^{d}.
\]
This proves coercivity. In particular, $Q_t$ attains a global minimum.
The origin is a critical point because $d\geq4$.

If $t>m_n$, Proposition~\ref{prop:threshold} gives $v$ with $P_t(v)<0$.
For sufficiently small $s>0$,
\[
 Q_t(sv)=s^d\left(P_t(v)+s^2\sum_a v_a^{d+2}\right)<0.
\]
A global minimizer of $Q_t$ is then a nonzero real critical point.

Conversely, at a nonzero real critical point $v$, Euler's identity gives
\begin{equation}\label{eq:euler}
 0=v\cdot\nabla Q_t(v)
   =dP_t(v)+(d+2)\sum_a v_a^{d+2}.
\end{equation}
The last sum is positive, so $P_t(v)<0$. By
Proposition~\ref{prop:threshold}, this implies $t>m_n$.
\end{proof}

\begin{remark}
At a nonzero real critical point, all gap coordinates are nonzero.
Indeed, $P_t(v)<0$ is incompatible with the nonnegativity of $N$ on
$D=0$. Thus the critical-point threshold also respects every collision
face of the node region.
\end{remark}

\section{A fixed quotient basis and its trace matrix}\label{sec:quotient}

In the coefficient ring $\Q[t]$, write $P(u,t)=N(u)-tD(u)$ and set
\begin{equation}\label{eq:relations}
 h_i(u,t)=-\frac{1}{d+2}\frac{\partial P}{\partial u_i}(u,t),
 \qquad f_i(u,t)=u_i^{d+1}-h_i(u,t).
\end{equation}
Thus $f_i=(d+2)^{-1}\partial Q_t/\partial u_i$.
Each $h_i$ is homogeneous of $u$-degree $d-1$ and has coefficients
affine in $t$. For a multi-index, write $\abs{\alpha}=\sum_i\alpha_i$.
Fix the lexicographic order on the set
\begin{equation}\label{eq:basis}
 \B=\{u^\alpha:\alpha\in\{0,1,\ldots,d\}^k\},
 \qquad \abs{\B}=M=(d+1)^k.
\end{equation}
The reduction rule is
\begin{equation}\label{eq:rewrite}
 u_i^{d+1}\longmapsto h_i(u,t).
\end{equation}

\begin{lemma}\label{lem:basis}
Reduction by~\eqref{eq:rewrite} terminates and gives a unique normal form
in the span of $\B$. The residue classes of $\B$ form a basis of
\begin{equation}\label{eq:quotient}
 \mathcal A=\Q[t,u_1,\ldots,u_k]/(f_1,\ldots,f_k)
\end{equation}
as a $\Q[t]$-module. For every real $\tau$, the same monomials form a
basis of
\[
 \mathcal A_\tau=\R[u_1,\ldots,u_k]/(f_1(u,\tau),\ldots,f_k(u,\tau)),
\]
and normal forms commute with specialization $t=\tau$.
\end{lemma}

\begin{proof}
Every replacement reduces total $u$-degree by two, so reduction
terminates. We prove uniqueness by induction on the total degree
of a monomial. If two possible first replacements use the same index,
their results agree. If they use distinct indices $i$ and $j$, write the
monomial as $u_i^{d+1}u_j^{d+1}v$. The two replacements give
$h_i u_j^{d+1}v$ and $u_i^{d+1}h_jv$. In the first expression replace
the displayed $u_j^{d+1}$ factor, and in the second replace the displayed
$u_i^{d+1}$ factor. Both paths reach $h_i h_jv$.
Every intermediate monomial has smaller degree than the original one.
The induction hypothesis makes all subsequent normal forms equal.
Linearity extends the conclusion to polynomials.

Let $\NF$ denote the resulting $\Q[t]$-linear normal-form map.
For a monomial $v$, the defining replacement shows
$\NF(vf_i)=0$. Hence $\NF$ annihilates the ideal $(f_1,\ldots,f_k)$.
Each reduction changes a polynomial by an element of this ideal, so
$p-\NF(p)$ belongs to the ideal for every $p$. These two facts show
that $\B$ spans the quotient and that a linear combination of $\B$
belonging to the ideal must be zero.

This proof works over $\R$ after any specialization. Moreover, the
reduction uses only addition and multiplication of coefficients and
the fixed rational factor $(d+2)^{-1}$. Specialization therefore commutes
with the entire reduction.
\end{proof}

For $g\in\Q[t,u]$, let $M_g(t)$ be the matrix of multiplication by $g$
in the ordered basis $\B$. Its column indexed by $u^\beta$ consists of
the coefficients of $\NF(gu^\beta)$. Define
\begin{equation}\label{eq:H}
 H_n(t)_{\alpha,\beta}
 =\tr\bigl(M_{u^{\alpha+\beta}}(t)\bigr),
 \qquad u^\alpha,u^\beta\in\B.
\end{equation}
This is a symmetric matrix over $\Q[t]$. Lemma~\ref{lem:basis} identifies
$H_n(\tau)$ with the trace-form matrix of $\mathcal A_\tau$ for every
real $\tau$.

\section{The real trace identity, including multiplicities}\label{sec:trace}

We give the finite-dimensional trace argument to account explicitly for
multiple zeros. For a finite-dimensional real algebra $A$, the trace
form is the symmetric bilinear form
\[
 (a,b)\longmapsto\tr(L_{ab}),
\]
where $L_c$ denotes multiplication by $c$ on $A$.

\begin{proposition}[Hermite trace identity]\label{prop:hermite}
Let $A=\R[u_1,\ldots,u_k]/J$ be a nonzero finite-dimensional algebra.
The signature of its trace form is the number of distinct real common
zeros of $J$.
\end{proposition}

\begin{proof}
Let $A_\C=A\otimes_\R\C$. The multiplication operators on $A_\C$
commute. They are simultaneously upper triangular in a suitable basis.
For completeness, a finite commuting family over $\C$ has a common
eigenvector: restrict the family to an eigenspace of one member and
continue, decreasing dimension whenever necessary. Its common invariant
line gives a commuting family on the quotient. Induction on dimension
then gives simultaneous upper triangularization. Applying this to the
multiplication operators of a vector-space basis of $A_\C$ suffices.

Every diagonal position defines a unital $\C$-algebra homomorphism
$\chi:A_\C\to\C$, since diagonal entries add and multiply under upper
triangular matrix operations. The values $\chi(u_i)$ give a common zero
of $J$. Let $Z$ be the finite set of zeros obtained in this way, and
let $m_\zeta$ be the number of diagonal positions yielding $\zeta$.
Then $m_\zeta$ is a positive integer and
\begin{equation}\label{eq:trace-sum}
 \tr(L_g)=\sum_{\zeta\in Z}m_\zeta g(\zeta)
 \qquad(g\in A_\C).
\end{equation}

Every complex common zero is in $Z$. To prove this, suppose $w$ is
a common zero outside $Z$. There is a polynomial $q$ over $\C$
with $q(w)=1$ and $q(\zeta)=0$ for all $\zeta\in Z$. Such a polynomial
is obtained by multiplying, for each $\zeta$, a normalized coordinate
difference separating $w$ and $\zeta$. The matrix $L_q$ is then strictly
upper triangular, so $L_q^{\dim A}=0$. Applying this operator to $1$
gives $q^{\dim A}=0$ in $A_\C$, contradicting evaluation at $w$.

Complex conjugation permutes $Z$. It also preserves the multiplicities.
Indeed, choose a polynomial $q_\zeta$ that equals one at $\zeta$ and
zero at the other points of $Z$. Conjugation of coefficients gives
\[
 \tr(L_{\overline{q_\zeta}})
 =\overline{\tr(L_{q_\zeta})}=m_\zeta.
\]
By~\eqref{eq:trace-sum}, the left side is $m_{\bar\zeta}$.

Write $r$ for the number of real points of $Z$ and choose one point
from each of its $c$ nonreal conjugate pairs. Evaluation defines a
surjective real linear map
\begin{equation}\label{eq:evaluation}
 A\longrightarrow\R^r\times\C^c.
\end{equation}
To see surjectivity, first interpolate arbitrary values at the finite
set $Z$, assigning conjugate values at conjugate points. Polynomial
interpolation at finitely many distinct points follows from the same
separating coordinate products. Averaging the interpolating polynomial
with its coefficientwise conjugate gives real coefficients and
preserves the assigned values.

For real $a,b$, equation~\eqref{eq:trace-sum} now expresses the trace
form as the pullback under~\eqref{eq:evaluation} of
\begin{equation}\label{eq:trace-decomposition}
 \sum_{\zeta\text{ real}}m_\zeta a(\zeta)b(\zeta)
 +\sum_{\{\zeta,\bar\zeta\}}2m_\zeta
       \operatorname{Re}\bigl(a(\zeta)b(\zeta)\bigr).
\end{equation}
Each real summand has one positive direction. On $\C\cong\R^2$,
the form $2m\operatorname{Re}(zw)$ has matrix
$\diag(2m,-2m)$ and hence one positive and one negative direction.
Choose a real vector-space complement to the kernel of evaluation.
Surjectivity identifies the trace form on that complement with
\eqref{eq:trace-decomposition}; the kernel contributes zero directions.
The signature is therefore $r$.
\end{proof}

\begin{corollary}\label{cor:signature-threshold}
For every real $t$,
\begin{equation}\label{eq:signature-threshold}
 \sig H_n(t)
 =\#\{u\in\R^k:\nabla Q_t(u)=0\},
 \qquad
 \sig H_n(t)=1\ \Longleftrightarrow\ t\leq m_n.
\end{equation}
\end{corollary}

\begin{proof}
Lemma~\ref{lem:basis} gives a quotient algebra of dimension $M$ at every
parameter. Its defining equations are nonzero scalar multiples of the
partial derivatives of $Q_t$. Apply Proposition~\ref{prop:hermite} and
then Proposition~\ref{prop:critical}.
\end{proof}

\section{The boundary polynomial and the selected root}\label{sec:selection}

Expand the characteristic polynomial as
\begin{equation}\label{eq:C}
 C_n(t,\lambda)=\det(\lambda I_M-H_n(t))
              =\sum_{j=0}^{M}c_j(t)\lambda^j.
\end{equation}
Since $c_M=1$, the following definition is finite and unambiguous.

\begin{definition}\label{def:R}
Let $e$ be the smallest index for which $c_e(t)$ is a nonzero polynomial,
and put $R_n(t)=c_e(t)$.
\end{definition}

\begin{lemma}\label{lem:degree}
Every entry of $H_n(t)$ has degree at most $kd$ in $t$, and
\begin{equation}\label{eq:degree}
 \deg R_n\leq kdM.
\end{equation}
\end{lemma}

\begin{proof}
Consider a diagonal coefficient of multiplication by $u^{\alpha+\beta}$
on the basis element $u^\gamma$. Reduction starts at total degree
$\abs{\alpha}+\abs{\beta}+\abs{\gamma}$. Every step decreases this
degree by exactly two. A term contributing to the coefficient of
$u^\gamma$ therefore uses exactly
\[
 \frac{\abs{\alpha}+\abs{\beta}}{2}\leq kd
\]
steps. If this number is not an integer, there is no such term.
The coefficient introduced by each step has $t$-degree at most one.
Thus every contributing coefficient has $t$-degree at most $kd$.
Taking the trace preserves the bound.

Each coefficient of the determinant in~\eqref{eq:C} is a sum of
products of at most $M$ matrix entries. Its degree is at most $kdM$,
which proves~\eqref{eq:degree}.
\end{proof}

\begin{lemma}\label{lem:boundary}
The signature of $H_n(t)$ is constant on each interval containing no
real root of $R_n$, and $R_n(m_n)=0$.
\end{lemma}

\begin{proof}
If $R_n(t)\ne0$, equation~\eqref{eq:C} shows that zero is an eigenvalue
of algebraic multiplicity exactly $e$. Symmetry identifies this
multiplicity with the nullity. The eigenvalues of a continuously varying
real symmetric matrix vary continuously. Near a given parameter where
$R_n\ne0$, the nonzero eigenvalues retain their signs, and the nullity
remains $e$. The signature is consequently locally constant there,
and hence constant on every such interval.

If $R_n(m_n)\ne0$, there would be an interval about $m_n$ on which the
signature is constant. Corollary~\ref{cor:signature-threshold} gives
signature one to the left of $m_n$ and a strictly larger signature to
its right, a contradiction.
\end{proof}

\begin{proof}[Proof of Theorem~\ref{thm:main}]
Attainment and the lower bound follow from Lemma~\ref{lem:attainment}.
The one-node integral is $2$. Lemma~\ref{lem:strict-upper} and
Proposition~\ref{prop:endpoints} prove the remaining assertions about
minimizers.

The explicit construction of $H_n$ is~\eqref{eq:H}, using the rational
forms~\eqref{eq:ND} and the reduction~\eqref{eq:rewrite}.
Definition~\ref{def:R} specifies $R_n$, and Lemma~\ref{lem:degree}
gives its degree bound. Since $0<m_n<2$, Lemma~\ref{lem:boundary}
shows that $m_n=r_h$ for some $1\leq h\leq s$.

For $j<h$, the rational sample $\rho_j$ is less than $m_n$; for
$j\geq h$, it is greater than $m_n$. By
Corollary~\ref{cor:signature-threshold},
\[
 \sig H_n(\rho_j)=1\quad\Longleftrightarrow\quad j<h.
\]
Exactly $h-1$ samples therefore have signature one, proving
\eqref{eq:main-rule}.

Finally, all polynomials and matrices in the construction are obtained
by finite rational arithmetic. The index $e$ is determined by testing
the finitely many rational coefficient polynomials for equality to
zero. Univariate real root isolation orders the roots of $R_n$ in
$(0,2)$ and supplies rational samples between them. The signatures
of the rational symmetric matrices at those samples are obtained by
exact symmetric elimination. These operations terminate for every
positive integer $n$.
\end{proof}

\begin{corollary}\label{cor:algebraic}
For $n\geq2$, the number $m_n$ is algebraic over $\Q$ of degree at most
\[
 2n(n-1)^2\bigl(2n(n-1)+1\bigr)^{n-1}.
\]
\end{corollary}

\begin{proof}
The minimal polynomial of $m_n$ divides the nonzero polynomial $R_n$.
Apply~\eqref{eq:degree}.
\end{proof}

\begin{remark}\label{rem:inertia}
Exact symmetric elimination requires only rational operations.
At a nonzero diagonal entry, a congruence and its Schur complement
remove a one-dimensional block. If all diagonal entries vanish and
some off-diagonal entry $b$ is nonzero, use the block
$\bigl(\begin{smallmatrix}0&b\\b&0\end{smallmatrix}\bigr)$, which has
one positive and one negative eigenvalue, and remove it by a block
Schur complement. An all-zero remaining matrix contributes only
zero eigenvalues. Dimension decreases at every step, including at
singular parameter values.
\end{remark}

\section{The two- and three-node cases}\label{sec:examples}

The first two nontrivial cases illustrate both the encoding and the
selected-root rule.

For $n=2$, one has $k=1$, $d=4$, and
\[
 D(u)=4u^4,\qquad N(u)=\frac{16}{3}u^4.
\]
Put $a=\frac83(t-\frac43)$. The normalized critical equation is
$u^5-au^3=0$, so reduction uses $u^5=au^3$. On the basis
$1,u,u^2,u^3,u^4$, the trace matrix is
\begin{equation}\label{eq:two-node-H}
 H_2(t)=
 \begin{pmatrix}
 5&0&2a&0&2a^2\\
 0&2a&0&2a^2&0\\
 2a&0&2a^2&0&2a^3\\
 0&2a^2&0&2a^3&0\\
 2a^2&0&2a^3&0&2a^4
 \end{pmatrix}.
\end{equation}
The trace moments follow directly from the multiplication matrix:
$\tr(M_1)=5$, $\tr(M_{u^{2j}})=2a^j$ for $j\geq1$, and all odd
moments vanish. Expanding the characteristic polynomial gives
\[
 R_2(t)=-12a^3(1+a^2)^2.
\]
Its only real root is $t=4/3$, and~\eqref{eq:main-rule} yields
$m_2=4/3$.

For $n=3$, Proposition~\ref{prop:endpoints} reduces the configuration
to $(-1,y,1)$ with $-1<y<1$. Its cardinal polynomials are
\[
 \frac{(x-y)(x-1)}{2(1+y)},\qquad
 \frac{1-x^2}{1-y^2},\qquad
 \frac{(x+1)(x-y)}{2(1-y)}.
\]
Their squared integrals sum to
\begin{equation}\label{eq:three-node}
 I(-1,y,1)=\frac43-\frac{4}{3(1-y^2)}
                   +\frac{8(3+y^2)}{15(1-y^2)^2}.
\end{equation}
As a function of $s=y^2\in[0,1)$, its derivative is
\[
 \frac{36+28s}{15(1-s)^3}>0.
\]
Thus the unique ordered minimizing triple is $(-1,0,1)$ and
$m_3=8/5$. Theorem~\ref{thm:main} selects this same value from $R_3$.

\section*{Acknowledgments}
The authors acknowledge assistance from AI systems developed by OpenAI
in mathematical exploration, checking, and preparation of this manuscript.

\begin{thebibliography}{9}

\bibitem{BPR}
S. Basu, R. Pollack, and M.-F. Roy,
\textit{Algorithms in Real Algebraic Geometry}, second ed.,
Algorithms and Computation in Mathematics, vol.~10,
Springer, Berlin, 2006.
\href{https://doi.org/10.1007/3-540-33099-2}{doi:10.1007/3-540-33099-2}.

\bibitem{Bloom}
T. F. Bloom (ed.), Erd\H{o}s Problem~1131,
\textit{Erd\H{o}s Problems}, online problem collection,
accessed November~11, 2026.
\href{https://www.erdosproblems.com/1131}{\nolinkurl{erdosproblems.com/1131}}.

\bibitem{BT}
L. Brutman and D. Toledano,
An extremal problem of Erd\H{o}s in interpolation theory,
\textit{Computers \& Mathematics with Applications} \textbf{34}
(1997), no.~12, 37--47.
\href{https://doi.org/10.1016/S0898-1221(97)00232-0}{doi:10.1016/S0898-1221(97)00232-0}.

\bibitem{ESVV}
P. Erd\H{o}s, J. Szabados, A. K. Varma, and P. V\'ertesi,
On an interpolation theoretical extremal problem,
\textit{Studia Scientiarum Mathematicarum Hungarica} \textbf{29}
(1994), 55--60.

\bibitem{LeSafey}
H. P. Le and M. Safey El Din,
Solving parametric systems of polynomial equations over the reals
through Hermite matrices,
\textit{Journal of Symbolic Computation} \textbf{112} (2022), 25--61.
\href{https://doi.org/10.1016/j.jsc.2021.12.002}{doi:10.1016/j.jsc.2021.12.002};
\href{https://arxiv.org/abs/2011.14136}{arXiv:2011.14136}.

\bibitem{White}
P. White, Erd\H{o}s \#1131: the four-node minimum, exactly,
\textit{Erd\H{o}s Problem a Day}, August~11, 2026,
\href{https://erdosproblemaday.com/day/1131-exact-four-node-minimum}{online article}.

\end{thebibliography}
\end{document}
